\documentclass[12pt]{article}
\usepackage[utf8]{inputenc}
\usepackage[T1]{fontenc}
\usepackage{hyperref}
\title{Vitou Language Resources for K-12 Curriculum Development}
\author{Compiled by Claude Sonnet 5 (claude-sonnet-5) for Andr\'as Kornai}
\date{}
\begin{document}
\maketitle

\section{Introduction}
Vitou (ISO 639-3: vto; also called Fitou, Manirem or Takar) is a Papuan language of the Tor family in Sarmi Regency, Indonesia. The sources give about 600--800 speakers, in Takar village and nearby Nengke. Speaker counts differ across sources and years. Vitou is not to be confused with Betaf, which is also called Manirem. Almost nothing educational was found.

\section{Primary Educational Resources}
None found. No primers, textbooks or school programmes were seen in the search results.

\section{Linguistic and Reference Materials}
\begin{itemize}
\item Wikipedia article on the Vitou language: \url{https://en.wikipedia.org/wiki/Vitou_language}
\item New Guinea World entry on Fitou (Tor coast): \url{https://newguineaworld.linguistik.uzh.ch/families/northwest-new-guinea/foja-range/orya-tor-river/tor-river/tor-coast/east-tor-coast/fitou}
\item Betaf language article (the other language called Manirem, for disambiguation): \url{https://en.wikipedia.org/wiki/Betaf_language}
\end{itemize}

\section{Religious and Community Resources}
\begin{itemize}
\item Joshua Project profile of Vitou in Indonesia (demographic information only): \url{https://joshuaproject.net/people_groups/18956}
\end{itemize}

\section{Digital and Audio Resources}
None found.

\section{Curriculum Development Context}
Vitou is spoken by older adults and appears to be moribund. Documentation and community-led work would have to come before any curriculum.
\end{document}